% ============================================================
%  EXAMEN TEMA 6: ESTEQUIOMETRÍA
%  Curso 4to Año — Química
%  Duración: 90 minutos
%  Temas: El mol y la masa molar, balanceo de ecuaciones y
%         cálculos estequiométricos
%  Autor: Ing. Gabriel Astudillo
%  Nota: enunciado limpio (sin pistas ni desarrollo). Las
%  soluciones están en la "Clave de Respuestas" al final.
% ============================================================

\documentclass[12pt, letterpaper]{article}

% ── Paquetes ─────────────────────────────────────────────────

\usepackage{mathwizards-guia}
\mwguia{Examen — Tema 6: Estequiometría}{Ing. Gabriel Astudillo $\cdot$ 4to Año}

\begin{document}
% ============================================================

% ── Portada / Título ─────────────────────────────────────────
\mwportada{Examen del Tema 6}{Estequiometría}{El Mol, el Balanceo de Ecuaciones y los Cálculos de las Reacciones}

\vspace{4pt}
\begin{cuadroinfo}
  \small
  \textbf{Nombre:} \rule{0.55\linewidth}{0.4pt} \quad \textbf{Fecha:} \rule{0.15\linewidth}{0.4pt} \\[6pt]
  \textbf{Tiempo disponible:} 90 minutos \quad$\bullet$\quad \textbf{Puntaje total:} 100 puntos \\[4pt]
  \textbf{Instrucciones:}
  \begin{enumerate}[leftmargin=*]
    \item Lee cada ejercicio con calma antes de resolverlo.
    \item Muestra \emph{todos} tus pasos; usa factores de conversión y cancela unidades.
    \item Al balancear, ajusta \emph{coeficientes} (nunca subíndices) y verifica al final.
    \item Datos: H $= 1$, C $= 12$, N $= 14$, O $= 16$, Na $= 23$, Al $= 27$, S $= 32$, Cl $= 35{,}5$, Ca $= 40$.
    \item No se permite el uso de calculadora.
    \item Escribe con letra clara y revisa tu examen antes de entregarlo.
  \end{enumerate}
\end{cuadroinfo}

\vspace{8pt}

% ============================================================
\section{El Mol y la Masa Molar (32 puntos)}
% ============================================================

\begin{enumerate}[leftmargin=*]
  \item \textbf{(10 pts)} Resuelve:
        \begin{enumerate}[label=\alph*)]
          \item \textquestiondown Cuántos moles hay en $54\;\text{g}$ de agua, $\text{H}_2\text{O}$? ($M = 18\;\text{g/mol}$)
          \item \textquestiondown Cuántos gramos pesan $2{,}5\;\text{mol}$ de $\text{CO}_2$? ($M = 44\;\text{g/mol}$)
          \item Calcula la masa molar del ácido sulfúrico, $\text{H}_2\text{SO}_4$.
        \end{enumerate}

  \item \textbf{(8 pts)} \textquestiondown Cuántas moléculas hay en $0{,}5\;\text{mol}$ de metano, $\text{CH}_4$? ($N_A = 6{,}022 \times 10^{23}$.)

  \item \textbf{(8 pts)}
        \begin{enumerate}[label=\alph*)]
          \item \textquestiondown Cuántos moles hay en $40\;\text{g}$ de hidróxido de sodio, $\text{NaOH}$? ($M = 40\;\text{g/mol}$)
          \item \textquestiondown Cuántos gramos pesan $4\;\text{mol}$ de cloruro de sodio, $\text{NaCl}$? ($M = 58{,}5\;\text{g/mol}$)
        \end{enumerate}

  \item \textbf{(6 pts)} \textquestiondown Cuántos átomos hay en $2\;\text{mol}$ de sodio? ($N_A = 6{,}022 \times 10^{23}$.)
\end{enumerate}

\newpage

% ============================================================
\section{Balanceo de Ecuaciones Químicas (34 puntos)}
% ============================================================

\begin{enumerate}[leftmargin=*]
  \item \textbf{(9 pts)} Balancea las siguientes ecuaciones:
        \[
          a)\ \text{N}_2 + \text{H}_2 \rightarrow \text{NH}_3 \qquad b)\ \text{Mg} + \text{O}_2 \rightarrow \text{MgO}
        \]

  \item \textbf{(9 pts)} Balancea la combustión del propano:
        \[
          \text{C}_3\text{H}_8 + \text{O}_2 \rightarrow \text{CO}_2 + \text{H}_2\text{O}
        \]

  \item \textbf{(8 pts)} En la ecuación balanceada $2\,\text{H}_2 + \text{O}_2 \rightarrow 2\,\text{H}_2\text{O}$: \textquestiondown cuántos moles de $\text{O}_2$ reaccionan con $6\;\text{mol}$ de $\text{H}_2$?

  \item \textbf{(8 pts)}
        \begin{enumerate}[label=\alph*)]
          \item Enuncia la ley de conservación de la masa.
          \item \textquestiondown Qué diferencia hay entre un coeficiente y un subíndice? \textquestiondown Cuál se puede cambiar al balancear?
        \end{enumerate}
\end{enumerate}

\newpage

% ============================================================
\section{Cálculos Estequiométricos, Composición y Fórmulas (34 puntos)}
% ============================================================

\begin{enumerate}[leftmargin=*]
  \item \textbf{(10 pts)} En la reacción $\text{N}_2 + 3\,\text{H}_2 \rightarrow 2\,\text{NH}_3$: \textquestiondown cuántos gramos de amoniaco ($\text{NH}_3$) se producen a partir de $6\;\text{g}$ de hidrógeno?

  \item \textbf{(8 pts)} Calcula la composición porcentual del dióxido de carbono, $\text{CO}_2$ ($M = 44\;\text{g/mol}$).

  \item \textbf{(8 pts)} Un compuesto tiene fórmula empírica $\text{CH}_2$ y masa molar $28\;\text{g/mol}$. \textquestiondown Cuál es su fórmula molecular?

  \item \textbf{(8 pts)} En la reacción $2\,\text{Al} + 3\,\text{Cl}_2 \rightarrow 2\,\text{AlCl}_3$: \textquestiondown cuántos gramos de $\text{AlCl}_3$ se producen con $27\;\text{g}$ de aluminio?
\end{enumerate}

\newpage

% ============================================================
\ifmwclaves
  \section{Clave de Respuestas}
  % ============================================================

  \subsection*{Sección 1: El Mol y la Masa Molar}

  \begin{enumerate}[leftmargin=*, label=\textbf{\arabic*.}]
    \item a) $n = \dfrac{54}{18} = 3\;\text{mol de H}_2\text{O}$.
          \quad b) $m = 2{,}5 \times 44 = 110\;\text{g de CO}_2$.
          \quad c) $M(\text{H}_2\text{SO}_4) = 2(1) + 32 + 4(16) = 2 + 32 + 64 = 98\;\text{g/mol}$.

    \item $0{,}5 \times 6{,}022 \times 10^{23} = 3{,}011 \times 10^{23}$ moléculas de $\text{CH}_4$.

    \item a) $n = \dfrac{40}{40} = 1\;\text{mol de NaOH}$.
          \quad b) $m = 4 \times 58{,}5 = 234\;\text{g de NaCl}$.

    \item $2 \times 6{,}022 \times 10^{23} = 1{,}2044 \times 10^{24}$ átomos de Na.
  \end{enumerate}

  \newpage

  \subsection*{Sección 2: Balanceo de Ecuaciones Químicas}

  \begin{enumerate}[leftmargin=*, label=\textbf{\arabic*.}]
    \item a) $\text{N}_2 + 3\,\text{H}_2 \rightarrow 2\,\text{NH}_3$ (N: 2 = 2; H: 6 = 6). \quad
          b) $2\,\text{Mg} + \text{O}_2 \rightarrow 2\,\text{MgO}$ (Mg: 2 = 2; O: 2 = 2).

    \item $\text{C}_3\text{H}_8 + 5\,\text{O}_2 \rightarrow 3\,\text{CO}_2 + 4\,\text{H}_2\text{O}$.
          Verificación: C: 3 = 3; H: 8 = 8; O: 10 = 6 + 4. \checkmark

    \item La proporción es $2$ mol $\text{H}_2$ : $1$ mol $\text{O}_2$. Con $6$ mol de $\text{H}_2$ reaccionan $3$ mol de $\text{O}_2$.

    \item a) La masa total de los reactivos es igual a la masa total de los productos (la materia no se crea ni se destruye, solo se transforma).
          \quad b) El \textbf{coeficiente} indica cuántas moléculas o moles participan (número grande, delante de la fórmula); el \textbf{subíndice} indica cuántos átomos de cada elemento hay dentro de la molécula. Al balancear solo se cambian los \textbf{coeficientes}.
  \end{enumerate}

  \newpage

  \subsection*{Sección 3: Cálculos Estequiométricos, Composición y Fórmulas}

  \begin{enumerate}[leftmargin=*, label=\textbf{\arabic*.}]
    \item Moles de $\text{H}_2$: $n = \dfrac{6}{2} = 3$ mol.
          Proporción $3$ mol $\text{H}_2$ : $2$ mol $\text{NH}_3$:
          \[
            3 \text{ mol H}_2 \cdot \frac{2 \text{ mol NH}_3}{3 \text{ mol H}_2} = 2 \text{ mol NH}_3.
          \]
          Masa de $\text{NH}_3$: $M(\text{NH}_3) = 14 + 3(1) = 17\;\text{g/mol}$, así que $m = 2 \times 17 = 34\;\text{g de NH}_3$.

    \item $M(\text{CO}_2) = 12 + 2(16) = 44\;\text{g/mol}$.
          \[
            \% \text{C} = \frac{12}{44} \times 100 \approx 27{,}3\%, \qquad \% \text{O} = \frac{32}{44} \times 100 \approx 72{,}7\%.
          \]

    \item $M_{\text{empírica}}(\text{CH}_2) = 12 + 2(1) = 14\;\text{g/mol}$. \quad $k = \dfrac{28}{14} = 2$.
          Fórmula molecular: $\text{C}_2\text{H}_4$ (eteno).

    \item Moles de Al: $n = \dfrac{27}{27} = 1$ mol. Proporción $2$ mol Al : $2$ mol $\text{AlCl}_3$ (es decir, 1 : 1): se producen $1$ mol de $\text{AlCl}_3$.
          \[
            M(\text{AlCl}_3) = 27 + 3(35{,}5) = 27 + 106{,}5 = 133{,}5\;\text{g/mol} \Rightarrow m = 1 \times 133{,}5 = 133{,}5\;\text{g}.
          \]
  \end{enumerate}

\fi
\end{document}
